\section{Introdução}

\begin{frame}{A Lei da Troca Equivalente}
  \begin{quote}
    ``Para se obter algo, é preciso oferecer algo em troca de igual valor.''\\
    \raggedleft \small -- Hiromu Arakawa, \textit{Fullmetal Alchemist}
  \end{quote}
  
  \vspace{1.5em}
  
  \textbf{A Metáfora da Dívida Técnica (DT):}
  \begin{itemize}
    \item \textbf{O ganho imediato:} Velocidade de entrega inicial, pragmatismo e validação rápida.
    \item \textbf{A troca (os juros):} Acúmulo de complexidade, defasagem documental e severo comprometimento na etapa de testes.
  \end{itemize}
\end{frame}

\begin{frame}{A Engenharia de Software e a DT}
  \textbf{Mecanismos de Prevenção:}
  \begin{itemize}
      \item Documentação formal do sistema (como documentos de requisitos).
      \item Validação dos requisitos do sistema.
      \item Verificação do sistema implementado diante dos requisitos.
      \item Análise de cobertura dos casos de teste.
  \end{itemize}

  \vspace{1em}

  \textbf{Por que não são executados?}
  \begin{itemize}
    \item \textbf{O Custo:} O elevado esforço inerente à sua manutenção, expresso em tempo e recursos financeiros.
    \item \textbf{Decisão de Curto Prazo:} A equipe é induzida a priorizar ganhos imediatos de codificação, mesmo ciente dos prejuízos futuros.
    \item \textbf{O Resultado:} Incorre-se formalmente na contração de DTs.
  \end{itemize}
  
\end{frame}

\begin{frame}{Objetivos}
  \textbf{Objetivo Geral:}
  \begin{itemize}
    \item Relatar a análise e reestruturação da documentação de requisitos e testes no Sistema de Fomento ao Microempreendedor (SFM), diagnosticando os impactos da DT na equipe.
  \end{itemize}
  
  \vspace{1em}
  
  \textbf{Objetivos Específicos:}
  \begin{itemize}
    \item Formalizar um documento de requisitos.
    \item Mapear a rastreabilidade e avaliar a cobertura de testes.
    \item Investigar a percepção da equipe técnica sobre a evolução do sistema.
    \item Mitigar os ``juros'' acumulados pela ausência de processos formais.
  \end{itemize}
\end{frame}